“穿越沙漠”是一款生存型游戏，需要选择最佳游戏策略，使玩家在保证到达终点的情况下尽可能获得更多资金。
%
针对“穿越沙漠”问题中不同信息条件下的路径规划、资源配置与多人协同行动决策，本文以剩余资金最大化为核心目标，
%
建立了单人确定性决策模型、单人随机决策模型以及多人协同优化模型，对不同关卡情形下的最优策略进行了系统研究。
%
首先，根据地图中各区域的邻接关系构建无向图模型，利用Dijkstra 算法求解起点、村庄、矿山与终点之间的最短路径，
%
为后续策略优化提供路径基础。

对于问题一，我们认为玩家不可能在任意结点停留，因此我们首先证明：除了沙暴天气必须停留以外，玩家只可能在矿山
%
停留挖矿，其余地点不做停留；然后，我们再根据游戏规则，给出玩家向前行走、沙暴时原地停留、矿山挖矿、村庄补给
%
等情况下生活必需品以及资金剩余量的递推关系式的基础之上，以剩余资金最多为目标，设置了玩家按时到达、生活必需
%
品重量不超过负重上限的约束条件，建立了单目标最优化模型，采用动态规划和回溯法，使用$C\mathbin{++}$编程求解。

对于问题二，在仅知道当天天气、未来天气具有随机性的条件下，将天气演化过程引入马尔可夫决策思想，
%
建立基于状态转移概率的动态决策模型。模型结合当前位置、剩余资源、剩余天数及资金状态，
%
给出玩家向前行走、沙暴时原地停留、矿山挖矿、村庄补给等动作进行滚动优化，从而获得不确定天气
%
条件下的最优策略，提高了模型在动态环境中的适应能力。

对于问题三，考虑多人同时参与时因同行、同购和同矿带来的资源消耗放大与收益分摊效应，
%
我们选择合作而非竞争的的博弈模式，在单人模型基础上建立多人协同优化框架。
%
进一步采用蒙特卡洛方法模拟天气过程，对不同随机天气样本下的多人协同行为进行数值实验，
%
从而评估策略的稳定性与收益表现。结果表明，所建模型能够较好统筹路径选择、资源补给与行动安排，
%
在确定与不确定天气条件下均具有较强的适用性与可解释性。

本文提出的模型将最短路径搜索、动态规划、马尔可夫决策与蒙特卡洛模拟有机结合，
%
可为复杂环境下的资源受限路径优化与群体协同决策问题提供参考。